\input{sol_pre.tex}
\input{q_2026-03-02_1.tex}
\begin{document}
\input{sol_head.tex}

%================= 1
\Savol{1}
12 ga böliniş $\iff$ 3 ga va 4 ga böliniş.\\
3 ga: $8+9+6+2+a+b=25+a+b$ $\Rightarrow$ $a+b\equiv2\pmod3$. \ 4 ga: $\overline{ab}$ 4 ga bölinadi.\\
Mos $\overline{ab}$: $08,20,32,44,56,68,80,92$. \ Köpaytmalar: $0,0,6,16,30,48,0,18$. Eng kattasi $6\cdot8=48$.
\Javob{B) 48}

%================= 2
\Savol{2}
Surat: $45\frac{10}{63}-44\frac{25}{84}=1+\frac{40-75}{252}=1-\frac5{36}=\frac{31}{36}$.\\
Maxraj: $\left(\frac73-\frac{10}9\right):4-\frac34=\frac{11}{36}-\frac{27}{36}=-\frac{16}{36}$.
\[\frac{31}{36}:\left(-\frac{16}{36}\right):\frac{31}{32}=-\frac{31}{16}\cdot\frac{32}{31}=-2.\]
\Javob{B) $-2$}

%================= 3
\Savol{3}
\[0{,}1x+\frac y6=\frac x2-\frac y2\ \Rightarrow\ \frac23y=0{,}4x\ \Rightarrow\ \frac xy=\frac{2/3}{0{,}4}=\frac53.\]
\Javob{B) $\frac53$}

%================= 4
\Savol{4}
$x+y=205\,000$, \ $0{,}85x=1{,}2y\Rightarrow x=\frac{24}{17}y$ $\Rightarrow$ $y=85\,000$, \ $x=120\,000$. \ Farq: $35\,000$.
\Javob{B) 35000}

%================= 5
\Savol{5}
$-1^k=-1$, şuning uçun hadlar $(-1)^{k+1}$, $k=2,\dots,99$: \ $-1+1-1+\dots+1$ (98 ta had, juft-juft nol).
\Javob{A) 0}

%================= 6
\Savol{6}
Kvadratlarini soliştiramiz ($\sqrt{96}=4\sqrt6=2\sqrt{24}$):
\[a^2=11+\sqrt{96},\qquad b^2=10+\sqrt{96},\qquad c^2=14+\sqrt{96}\ \Rightarrow\ c>a>b.\]
\Javob{D) $c>a>b$}

%================= 7
\Savol{7}
$(2\pm\sqrt3)^3=8\pm12\sqrt3+18\pm3\sqrt3=26\pm15\sqrt3$:
\[(2+\sqrt3)+(2-\sqrt3)=4.\]
\Javob{C) 4}

%================= 8
\Savol{8}
$\dfrac{S_n}n=2n\Rightarrow S_n=2n^2$, \ $a_{10}=S_{10}-S_9=200-162=38$.
\Javob{A) 38}

%================= 9
\Savol{9}
$b_2^3=8000\Rightarrow b_2=20$; \ $b_1+b_3=50$, $b_1b_3=400$ $\Rightarrow$ $\{10;40\}$, \ $q=2$ yoki $q=\frac12$.\\
$q=2$: $10+20+40+80+160=310$; \quad $q=\frac12$: $40+20+10+5+2{,}5=77{,}5$.
\Javob{A) 310 yoki 77,5}

%================= 10
\Savol{10}
Surat: $\dfrac{x^2+2xy-4x^2+y^2}x=\dfrac{y^2+2xy-3x^2}x=\dfrac{(y+3x)(y-x)}x$. \ Maxraj: $y(y^2+2xy-3x^2)=y(y+3x)(y-x)$.
\[\frac{(y+3x)(y-x)}{x\cdot y(y+3x)(y-x)}=\frac1{xy}.\]
\Javob{D) $\frac1{xy}$}

%================= 11
\Savol{11}
\[\frac12\left(\frac{x^2-2}{2x}\right)^2+1=\frac{x^4-4x^2+4+8x^2}{8x^2}=\frac{(x^2+2)^2}{8x^2}\ \Rightarrow\ \sqrt{\ }=\frac{x^2+2}{2\sqrt2\,x}.\]
\[\frac{x^2+2}{2\sqrt2\,x}:\frac{x^2+2}{2x}=\frac1{\sqrt2}.\]
\Javob{B) $\frac1{\sqrt2}$}

%================= 12
\Savol{12}
IV çorakda $\cos\alpha>0$: \ $\cos\alpha=\sqrt{1-\frac{16}{25}}=\frac35$.
\Javob{A) $\frac35$}

%================= 13
\Savol{13}
$\sqrt3\sin^2x=\frac{\sqrt3}2-\frac{\sqrt3}2\cos2x$, \ $\sin^22x=1-\cos^22x$. \ $c=\cos2x$:
\[1-c^2+\frac{\sqrt3}2c+\frac12=0\ \Rightarrow\ c^2-\frac{\sqrt3}2c-\frac32=0\ \Rightarrow\ c=-\frac{\sqrt3}2\ \ (c=\sqrt3>1).\]
$2x=\pm\frac{5\pi}6+2\pi k\Rightarrow x=\pm\frac{5\pi}{12}+\pi k$. \ $\left[-\frac\pi2;\pi\right]$ da: $-\frac{5\pi}{12}$, $\frac{5\pi}{12}$, $\frac{7\pi}{12}$ — 3 ta.
\Javob{C) 3}

%================= 14
\Savol{14}
$2^{-(x^2+3x)}>2^{-x-3}\Rightarrow -x^2-3x>-x-3\Rightarrow x^2+2x-3<0\Rightarrow -3<x<1$.\\
Butunlar: $-2,-1,0$ — 3 ta.
\Javob{B) 3}

%================= 15
\Savol{15}
$0<\log_2L<1\Rightarrow 1<L<2$, \ $L=\log_{x-1}9$. $L>0$ va $9>1$ $\Rightarrow$ asos $x-1>1$:
\[1<\log_{x-1}9<2\ \Rightarrow\ x-1<9<(x-1)^2\ \Rightarrow\ 3<x-1<9\ \Rightarrow\ 4<x<10.\]
Butunlar: $5,6,7,8,9$ — 5 ta.
\Javob{D) 5 ta}

%================= 16
\Savol{16}
$(x+6)(x-1)=x^2+5x-6$, \ $(x+2)(x+3)=x^2+5x+6$. \ $t=x^2+5x$:
\[\frac4{t-6}-\frac5{t+6}=\frac14\ \Rightarrow\ 16(t+6)-20(t-6)=t^2-36\ \Rightarrow\ t^2+4t-252=0\ \Rightarrow\ t=14\ \text{yoki}\ t=-18.\]
$x^2+5x-14=0$: $x=2;\,-7$. \quad $x^2+5x+18=0$: $D<0$. \ 2 ta ildiz.
\Javob{B) 2}

%================= 17
\Savol{17}
$t=x-7$: \ $(t+1)^4+(t-1)^4=16\Rightarrow 2t^4+12t^2+2=16\Rightarrow t^4+6t^2-7=0\Rightarrow t^2=1$.\\
$x=6$, $x=8$ — 2 ta.
\Javob{B) 2}

%================= 18
\Savol{18}
$x^2-2x+1=(x-1)^2$:
\[(x-1)\left(\frac{x+1}{x^2+1}+x-1\right)=\frac{(x-1)(x^3-x^2+2x)}{x^2+1}=\frac{x(x-1)(x^2-x+2)}{x^2+1}\le0.\]
$x^2-x+2>0$ $\Rightarrow$ $x(x-1)\le0\Rightarrow0\le x\le1$. Butunlar: $0;\ 1$.
\Javob{C) 2}

%================= 19
\Savol{19}
$(x-3)^2\ge0$: tengsizlik $x>0$, $x\ne3$ da bajariladi. Yeçim bölmaydigan natural son faqat $3$.
\Javob{A) 1}

%================= 20
\Savol{20}
$f(3)=\ln3+\ln3=\ln9$, \ $g(\ln9)=2e^{\ln9}=18$.
\Javob{C) 18}

%================= 21
\Savol{21}
$\dfrac{x^4}{x^2+1}=x^2-1+\dfrac1{x^2+1}$ (juft funksiya):
\[2\int_0^1\left(x^2-1+\frac1{x^2+1}\right)dx=2\left(\frac13-1+\frac\pi4\right)=\frac\pi2-\frac43.\]
\Javob{D) $\frac\pi2-\frac43$}

%================= 22
\Savol{22}
$y'=-\dfrac{2x}{(x^2+4)^2}=0\Rightarrow x=0$, \ $y(0)=\frac14$. \ Yiğindi: $0+\frac14=\frac14$.
\Javob{D) $\frac14$}

%================= 23
\Savol{23}
Kvadrat diagonali — diametr: $d=2\sqrt6$. \ $S=\dfrac{d^2}2=\dfrac{24}2=12$.
\Javob{B) 12}

%================= 24
\Savol{24}
$x^2=\dfrac{a^2-1}{(a-1)^2}=\dfrac{a+1}{a-1}$ $\Rightarrow$ $\dfrac{x^2+1}{x^2-1}=\dfrac{2a}2=a$.\\
$2y^2=\dfrac{1-b}{1+b}$ $\Rightarrow$ $\dfrac{2y^2-1}{2y^2+1}=\dfrac{-2b}2=-b$. \quad $4z^2=\dfrac{1-c^2}{(c-1)^2}=\dfrac{1+c}{1-c}$ $\Rightarrow$ $\dfrac{4z^2-1}{4z^2+1}=c$.
\[f=a-(-b)+c=a+b+c.\]
\Javob{A) $a+b+c$}

%================= 25
\Savol{25}
$3x-4$ va $2x+1$ butun $\Rightarrow$ $x=\frac{m+4}3$, $2x+1=\frac{2m+11}3$ butun $\Rightarrow$ $x$ butun.\\
Şartlar: $3x-4\le8$, $2x+1\le8$ ($8$ — eng katta), $(3x-4)+(2x+1)>8$ $\Rightarrow$ $2{,}2<x\le3{,}5$ $\Rightarrow$ $x=3$.\\
Tomonlar $5;\ 7;\ 8$, \ $P=20$.
\Javob{D) 20}

%================= 26
\Savol{26}
Uçinçi tomon (kosinuslar teoremasi):
\[c^2=1+(1+\sqrt3)^2+2\cdot1\cdot(1+\sqrt3)\cdot\frac12=6+3\sqrt3=3(2+\sqrt3)\ \Rightarrow\ c=\sqrt3\cdot\frac{\sqrt6+\sqrt2}2.\]
Sinuslar teoremasi:
\[\sin\alpha=\frac{1\cdot\sin120^\circ}c=\frac{\sqrt3/2}{\sqrt3(\sqrt6+\sqrt2)/2}=\frac1{\sqrt6+\sqrt2}=\frac{\sqrt6-\sqrt2}4=\sin15^\circ.\]
\Javob{C) $15^\circ$}

%================= 27
\Savol{27}
$KO=40$, $OL=8$. Teng yonli trapetsiyada $AH=\dfrac{AD+BC}2=KL=48$.\\
$OL$ — $\triangle CHD$ ning örta çiziği: $HD=2\cdot OL=16$. \quad $AD=48+16=64$.
\Javob{C) 64}

%================= 28
\Savol{28}
$25\le\dfrac{n(n-3)}2\le30\Rightarrow50\le n(n-3)\le60$. \ $n=9$: $9\cdot6=54$ — mos ($n=8$: 40, $n=10$: 70).
\Javob{A) 9}

%================= 29
\Savol{29}
$D$ — $AC$ örtasi: $D(2;1;1)$. \ $\overrightarrow{DB}=(0;1;0)$, \ $\overrightarrow{DC}=(1;2;2)$:
\[\cos\angle BDC=\frac{0+2+0}{1\cdot3}=\frac23.\]
\Javob{B) $\frac23$}

%================= 30
\Savol{30}
$HB=AH\ctg45^\circ=4$, \ $HC=AH\ctg60^\circ=\dfrac4{\sqrt3}$.
\[S_{HBC}=\frac12\cdot4\cdot\frac4{\sqrt3}\cdot\frac{\sqrt3}2=4,\qquad V=\frac13\cdot4\cdot4=5\frac13.\]
\Javob{C) $5\frac13$}

%================= 31
\Savol{31}
$A$ da $n$ ta element: \ $2^{n+1}-2^{n-1}=3\cdot2^{n-1}=12\Rightarrow n=3$. \ Qism töplamlar: $2^3=8$.
\Javob{B) 8}

%================= 32
\Savol{32}
$11k\le1000\Rightarrow k\le90$; \ $k$ 2 ga ham, 3 ga ham bölinmasin: \ $90-45-30+15=30$.
\Javob{D) 30}

%================= 33-35
\newpage
\Savol{33}
\fig{\begin{tikzpicture}[scale=0.3]
\draw[thick] (-6,0) rectangle (6,8);
\draw[thick] (-6,0)--(0,8)--(6,0);
\draw[thick,draw=okgreen] (0,3) circle (3);
\draw[ax] (0,0)--(0,8); \draw[axs] (0,3)--(2.4,4.8);
\node[left] at (-6,4){\scriptsize 8}; \node[below] at (3,0){\scriptsize 6}; \node[aux] at (4.2,4.6){\scriptsize 10};
\node[aux] at (0.8,3.3){\tiny$\rho$};
\end{tikzpicture}}
Konus asosi silindr asosi bilan bir xil ($r=6$), balandligi silindrniki $H$. \ $\pi\cdot6\cdot l=60\pi\Rightarrow l=10$, \ $H=\sqrt{100-36}=8$.

\textbf{33.} $S_{\text{sil}}=2\pi\cdot36+2\pi\cdot6\cdot8=168\pi$, \ $S_{\text{kon}}=36\pi+60\pi=96\pi$; \ farq $72\pi$.

\textbf{34.} $V_{\text{sil}}=\pi\cdot36\cdot8=288\pi$, \ $V_{\text{kon}}=\frac13\cdot288\pi=96\pi$; \ farq $192\pi$.

\textbf{35.} Öq kesimi — tomonlari $10,10,12$ uçburçak: $\rho=\dfrac S p=\dfrac{\frac12\cdot12\cdot8}{16}=3$, \ $V=\frac43\pi\cdot27=36\pi$.
\Javob{33 – D; \ 34 – E; \ 35 – A}

%================= 36
\newpage
\Savol{36}
$t=\sqrt[4]x\ge0$: \ $t^2+2t^4+2=(t^4+3)t\Rightarrow t^5-2t^4-t^2+3t-2=0$.\\
$t=2$ ildiz: \ $(t-2)\left(t^4-t+1\right)=0$.\\
$t\ge0$ da $t^4-t+1>0$: \ $t\le1$ da $1-t\ge0$ va $t^4\ge0$ (birga nol emas); $t>1$ da $t^4>t$.\\
Demak $t=2$, \ $x=16$.

a) 1 ta ildiz. \quad b) Yiğindi: $16$.
\Javob{a) 1; \quad b) 16}

%================= 37
\newpage
\Savol{37}
Şart: $\sin x\ne0$, $\cos x\ne0$, $\sin x\cos x\ge0$.
\[\sin^3x+\sin^2x\cos x+\cos^3x+\cos^2x\sin x=(\sin x+\cos x)(\sin^2x+\cos^2x)=\sin x+\cos x.\]
$\sin x+\cos x=2\sqrt{\sin x\cos x}$ $\Rightarrow$ $\sin x+\cos x\ge0$; kvadratga: $1+2\sin x\cos x=4\sin x\cos x\Rightarrow\sin2x=1$.\\
$x=\frac\pi4+\pi k$; \ $\sin x+\cos x\ge0$ faqat $x=\frac\pi4+2\pi k$ da.

a) $\dfrac\pi4$.

b) $\left[-\frac{5\pi}2;\frac{5\pi}2\right]$ da: $-\frac{7\pi}4$, $\frac\pi4$, $\frac{9\pi}4$ — 3 ta.
\Javob{a) $\frac\pi4$; \quad b) 3}

%================= 38
\newpage
\Savol{38}
a) $4x^4+1\ge2\cdot2x^2\cdot1=4x^2$ (örta arifmetik va örta geometrik) $\Rightarrow$ $f(x)\le1$, tenglik $x^2=\frac12$ da. \ Eng katta qiymat $1$.

b) $f(x)\ge0$, $f(0)=0$; uzluksizlikdan qiymatlar sohasi $[0;1]$. Butun sonlar: $0;\ 1$ — 2 ta.
\Javob{a) 1; \quad b) 2}

%================= 39
\newpage
\Savol{39}
a) $a>0$ da $\arctg a+\arctg\frac1a=\frac\pi2$. \ $f(x)+f(2026-x)=\arctg2^{x-1013}+\arctg2^{1013-x}=\frac\pi2$.\\
$0,\dots,2026$ — 2027 ta had: 1013 ta juftlik va örtadagi $f(1013)=\arctg1=\frac\pi4$:
\[1013\cdot\frac\pi2+\frac\pi4=\frac{2027\pi}4.\]
b) $f'(x)=\dfrac{2^{x-1013}\ln2}{1+4^{x-1013}}$, \ $f'(1013)=\dfrac{\ln2}{2}$.
\Javob{a) $\frac{2027\pi}4$; \quad b) $\frac{\ln2}2$}

%================= 40
\newpage
\Savol{40}
$\frac2x=2\Rightarrow x=1$. Böyalgan soha: $[0;1]$ da $y=2$ ostida, $[1;2]$ da $y=\frac2x$ ostida.

a) $S=2\cdot1+\displaystyle\int_1^2\frac2x\,dx=2+2\ln2$.

b) \[V=\pi\left(\int_0^1 4\,dx+\int_1^2\frac4{x^2}dx\right)=\pi\left(4+4\cdot\frac12\right)=6\pi\approx18.\]
\Javob{a) $2+2\ln2$; \quad b) 18}

%================= 41
\newpage
\Savol{41}
a) $AC=16$. Styuart teoremasi:
\[BD^2=\frac{AB^2\cdot DC+BC^2\cdot AD}{AC}-AD\cdot DC=\frac{600+1960}{16}-60=100\ \Rightarrow\ BD=10.\]
b) $\triangle ABD$ teng tomonli ($10,10,10$), $\angle ADB=60^\circ$, $\angle BDC=120^\circ$.\\
$DO_1$, $DO_2$ — qöşni burçaklar bissektrisalari $\Rightarrow$ $\angle O_1DO_2=90^\circ$.\\
$r_1=\dfrac{10}{2\sqrt3}$, \ $DO_1=\dfrac{r_1}{\sin30^\circ}=\dfrac{10}{\sqrt3}$. \quad $\triangle BDC$: $S=\frac12\cdot10\cdot6\cdot\sin120^\circ=15\sqrt3$, $p=15$, $r_2=\sqrt3$, \ $DO_2=\dfrac{r_2}{\sin60^\circ}=2$.
\[S_{O_1O_2D}=\frac12\cdot\frac{10}{\sqrt3}\cdot2=\frac{10\sqrt3}3.\]
\Javob{a) 10; \quad b) $\frac{10\sqrt3}3$}

%================= 42
\newpage
\Savol{42}
a) $S_{EFG}=\frac12\cdot8\cdot6=24$.

b) $\dfrac{BE}{BA}=\dfrac{BF}{BC}=\dfrac47$ $\Rightarrow$ $EF\parallel AC$, $EF=\frac47AC\Rightarrow AC=14$.\\
$\dfrac{CF}{CB}=\dfrac{CG}{CD}=\dfrac37$ $\Rightarrow$ $FG\parallel BD$, $FG=\frac37BD\Rightarrow BD=14$.\\
$EF\perp FG$ $\Rightarrow$ $AC\perp BD$:
\[S_{ABCD}=\frac12\cdot AC\cdot BD=\frac12\cdot14\cdot14=98.\]
\Javob{a) 24; \quad b) 98}

%================= 43
\newpage
\Savol{43}
Içki aylana: $AB+CD=3+5=8$; taşqi aylana $\Rightarrow$ teng yonli: $AB=CD=4$. \ Balandlik $h=\sqrt{16-1}=\sqrt{15}$, \ $r=\dfrac{\sqrt{15}}2$.\\
Urinmalar: $B$, $C$ dan $\frac32$; \ $A$, $D$ dan $\frac52$. $M$, $N$ — $BC$, $AD$ dagi uriniş nuqtalari.

a) $BCFO_1E$ = $BEO_1M$ + $MO_1FC$ (har biri ikki teng töğri burçakli uçburçak):
\[S=2\cdot\frac32\cdot r=3\cdot\frac{\sqrt{15}}2=\frac{3\sqrt{15}}2.\]
b) Xuddi şunday $S_{AEO_1FD}=2\cdot\frac52\cdot r=\frac{5\sqrt{15}}2$. Undan $\triangle AO_2D$ ni ayiramiz.\\
$O_2$ dan $AD$ gaça $y$: \ $R^2=\left(\frac52\right)^2+y^2=\left(\frac32\right)^2+(\sqrt{15}-y)^2\Rightarrow y=\dfrac{11}{2\sqrt{15}}$ \ ($y<r$, $O_2$ beşburçak içida).
\[S_{AO_2D}=\frac12\cdot5\cdot\frac{11}{2\sqrt{15}}=\frac{11\sqrt{15}}{12},\qquad S=\frac{5\sqrt{15}}2-\frac{11\sqrt{15}}{12}=\frac{19\sqrt{15}}{12}.\]
\Javob{a) $\frac{3\sqrt{15}}2$; \quad b) $\frac{19\sqrt{15}}{12}$}

%================= 44
\newpage
\Savol{44}
Yon qirralar teng $\Rightarrow$ uç asosga taşqi çizilgan aylana markazi $O$ ga proyeksiyalanadi.\\
Trapetsiya: balandlik $\sqrt{9-1{,}5^2}=\frac{3\sqrt3}2$, diagonal $d^2=3^2+5\cdot8=49\Rightarrow d=7$.\\
Taşqi aylana — tomonlari $3,7,8$ uçburçakniki: $S=\frac12\cdot8\cdot\frac{3\sqrt3}2=6\sqrt3$, \ $R=\dfrac{3\cdot7\cdot8}{4\cdot6\sqrt3}=\dfrac7{\sqrt3}$.

a) $H=\sqrt{l^2-R^2}=\sqrt{\dfrac{196}3-\dfrac{49}3}=\sqrt{49}=7$.

b) Şar markazi balandlikda; yon qirra va balandlikdan: $R_{sh}=\dfrac{l^2}{2H}=\dfrac{196/3}{14}=\dfrac{14}3$.
\Javob{a) 7; \quad b) $\frac{14}3$}

%================= 45
\newpage
\Savol{45}
$AD=\sqrt{50^2-30^2}=40$ m. \ $AB=\dfrac{30}{\sin\alpha}$, \ $BC=40-30\ctg\alpha$ (metrda).
\[t(\alpha)=\frac{30}{9\sin\alpha}+\frac{40-30\ctg\alpha}{15}\ \ \left(\frac{\text{m}}{\text{km/soat}}\right),\qquad t'(\alpha)=\frac{-30\cos\alpha}{9\sin^2\alpha}+\frac{2}{\sin^2\alpha}=0.\]
a) $\cos\alpha=\dfrac{2\cdot9}{30}=\dfrac35$.

b) $\sin\alpha=\frac45$, $\ctg\alpha=\frac34$: \ $AB=37{,}5$ m, \ $BC=40-22{,}5=17{,}5$ m.
\[t=\frac{0{,}0375}9+\frac{0{,}0175}{15}=\frac1{240}+\frac7{6000}=\frac{25+7}{6000}=\frac{2}{375}\ \text{soat}.\]
\Javob{a) $\frac35$; \quad b) $\frac2{375}$}

\ShowJavoblar
\end{document}
